\scanpage{83}
\mspage{89}

% slanted
\margin{Mais dans le cas du \unsure{rev.} à 2 feuillets de $S$, il n'y a pas même de calcul à faire !}

Voici \deleted{pourquoi} le \deleted{petit} calcul qui montre que $f:X\to S$ est une $\mathbb D$-équivalence. On fait le calcul en cohomologie, on \unsure{doit|désire} prouver que les

\[
H^{\unsure{\cdot|i}}(S,k)\longrightarrow H^{\cdot}(X,k)
\]

iso (et de $=$ pour tout coeff. constant \deleted{\unsure{sur}} $M$ sur $S$).

Or $H^{\cdot}(X,k)\overset{\mathrm{def}}=R\Gamma_X(k_X)=\unsure{\mathbb R|R}\Gamma_S(Rf_*(k_X))$, et ceci est induit par

\[
i:k_S\unsure{\relbar|\longrightarrow} Rf_*k_X
\]

Or on a

\[
Rf_*k_X\simeq\underbrace{k_S}_{\text{\unsure{im.}}(i)}\oplus
\underbrace{k_S[1]}_{k_S\ \text{en degré }1}^{\text{\uline{tordu}}}\,,
\]

savoir \textquotedblleft tordu\textquotedblright\ par l'action de monodromie de $S=S^1$, laquelle est la symétrie \deleted{\unsure{$-1$}} $\lambda\longmapsto-\lambda$. On a

\[
R\Gamma_S(k_S[1]^{\text{tordu}}\unsure{]|)}=
H^1(\underset{\substack{\parallel\\\mathbb Z}}{\underline{\pi_1(S)}},k^{\text{tordu}}\unsure{)|\}} )=
\underset{\scriptstyle k/(u-\mathrm{id})k}{k_{(\mathbb Z)}}=\underset{\substack{\parallel\\0}}{k/2k}
\]

d'où le résultat.

Notons que si on regarde

\[
X'=X\times_S X\xrightarrow{f'}S\ ,
\]

on aura par Künneth relatif

\[
\begin{aligned}
Rf'_*(k_{X'})&\simeq Rf_*(k_X)\mathbin{\overset{\unsure{\mathbb L|L}}{\otimes}}Rf_*(k_X)\simeq\\
&k_S\oplus k_S[1]^t\oplus k_S[1]^t+
\underbrace{k_S(1)^t\otimes k_S[1)^t}_{\displaystyle k_S[2]^{\text{\uline{n. tordu}}}}\, ,
\end{aligned}
\]

\deleted{\unsure{donc}} d'où

\[
\begin{aligned}
\unsure{R H|H}^{\cdot}_{\mathbb D}(X',k)&\simeq R\Gamma_{X'}(k_{X'})\simeq R\Gamma_S(Rf'_*(k_{X'})\\
&\simeq R\Gamma_S(k_S)\oplus R\Gamma_S(k_S)[2]
\end{aligned}
\]

et $H^{\cdot}_{\mathbb D}(\unsure{S|X},k)\to H^{\cdot}_{\mathbb D}(X',k)$ \unsure{n'est} \uline{pas} iso.

% archive sheet number: 82
